% Input after the outer weak-type applications defining W_\tau and B_\tau.
\subsection{小 \texorpdfstring{$\tau$}{tau} 的单原子下界}

继续记 $w_\tau(x)=|x|^{\tau-1}$，$0<\tau<1$，$W_\tau$ 为其外层弱型常数。
下面把经典幂函数系数 $M_cw_\tau(1)$ 用作单原子测试下界。
其标量最大化与唯一半径已见 Dror--Ganguli--Strichartz
\cite[Lemma~3.3, (3.8)]{DrorGanguliStrichartz1996}；这里给出当前加权弱型模型的实现和有理证书。
仍记前述双原子界 $B_\tau=3(5/3)^\tau/(4\tau)$。

\begin{theorem}[单原子测试与显式下界]
\label{thm:single-atom-small-tau}
对每个 $r>0$，有
\[
 W_\tau\ge J_\tau(r):=
 \begin{cases}
 \displaystyle\frac{(1+r)^\tau-(1-r)^\tau}{2\tau r},&0<r\le1,\\[6pt]
 \displaystyle\frac{(r-1)^\tau+(r+1)^\tau}{2\tau r},&r\ge1.
 \end{cases} \tag{TA.1}
\]
因而 $W_\tau\ge\max\{B_\tau,\sup_{r>0}J_\tau(r)\}$。当
$0<\tau<1$，$J_\tau$ 在单原子族上有唯一极大点 $r_\tau>1$，它由
\[
 \tau r\bigl[(r-1)^{\tau-1}+(r+1)^{\tau-1}\bigr]
 =(r-1)^\tau+(r+1)^\tau \tag{TA.2}
\]
唯一确定；这是单原子参数族的结论，不是 $W_\tau$ 的全局最优化结论。
特别地，
\[
 W_{1/8}\ge J_{1/8}(17/16)>\frac{27}{4}>B_{1/8},\qquad
 \liminf_{\tau\downarrow0}\tau W_\tau\ge1. \tag{TA.3}
\]
\end{theorem}

\begin{proof}
对任意位于 $a\ne0$ 的单原子 $m\delta_a$，反射和伸缩归一化为原子位置 $1$，参数为 $r=m/(2|a|)$。先把输入取为 $\sigma_r=2r\delta_1$，阈值取 $1$。严格原子光晕恰为
\[
 H_r=(1-r,1+r). \tag{TA.4}
\]
若区间含原子，其半径至少为 $|x-1|$；反之当 $|x-1|<r$ 可取半径严格介于
$|x-1|$ 与 $r$ 之间。端点处最佳平均恰等于 $1$，故被严格水平集排除。
将 (TA.4) 对 $|x|^{\tau-1}$ 积分并除以输入成本 $2r$，即得 (TA.1)。

原子输入严格转回允许的函数类：令
\[
 f_\varepsilon=\frac r\varepsilon\mathbf1_{(1-\varepsilon,1+\varepsilon)}.
\]
则 $\|f_\varepsilon\|_{L^1(\mu_\tau)}\to2r$。若 $x\in H_r$，选
$|x-1|<\rho<r$，再令 $\varepsilon<\rho-|x-1|$；则 $Q(x,\rho)$ 含整个小区间输入，
且其平均为 $r/\rho>1$。故 Fatou 引理给出 $W_\tau\ge J_\tau(r)$。

对 $r>1$，令 $N(r)=(r-1)^\tau+(r+1)^\tau$、$F(r)=rN'(r)-N(r)$。有
\[
 F'(r)=rN''(r)=\tau(\tau-1)r\bigl[(r-1)^{\tau-2}+(r+1)^{\tau-2}\bigr]<0.
\]
又 $F(r)\to+\infty$ 当 $r\downarrow1$，而
\(F(r)\sim2(\tau-1)r^\tau\to-\infty\) 当 $r\to\infty$。因
$J_\tau'(r)=F(r)/(2\tau r^2)$，(TA.2) 有唯一根且给出 $(1,\infty)$ 上的唯一极大点。
对 $0<r<1$，令 $L(r)=(1+r)^\tau-(1-r)^\tau$。此时
\[
 L''(r)=\tau(\tau-1)\bigl[(1+r)^{\tau-2}-(1-r)^{\tau-2}\bigr]>0.
\]
所以 $rL'(r)-L(r)$ 从零严格增加，$L(r)/r$ 严格增加到 $r=1$；而右支在
$1+$ 处先增，故全单原子极大点确在 $r_\tau>1$。

为证 (TA.3)，
\[
 J_{1/8}(17/16)=\frac{64}{17}\left(16^{-1/8}+(33/16)^{1/8}\right).
\]
整数算术验证
\[
 (707/1000)^2<1/2,\qquad (547/500)^8<33/16,\qquad (13/12)^8>5/3.
\]
从而
\[
 J_{1/8}(17/16)>\frac{64}{17}
 \left(\frac{707}{1000}+\frac{547}{500}\right)
 =\frac{14408}{2125}>\frac{27}{4}>\frac{13}{2}
 >6(5/3)^{1/8}=B_{1/8}.
\]
最后取 $r=1+\tau$，有
\[
 \tau J_\tau(1+\tau)=\frac{\tau^\tau+(2+\tau)^\tau}{2(1+\tau)}\longrightarrow1,
\]
故 $W_\tau\ge J_\tau(1+\tau)$ 给出所述 liminf。
\end{proof}

由 (OW.2) 的等指数双线性特例，对 $w(x)=|x|^{-7/8}$ 还得到
\[
 \lim_{\lambda\downarrow1/4}(\lambda-1/4)
 C^{(2),\mathrm{out}}_{c,w}(\lambda)
 =\tfrac12W_{1/8}>\frac{27}{8}. \tag{TA.5}
\]
